\chapter*{Nota}
El objetivo de hacer estos apuntes no es otro sino el de darles un texto de apoyo a las diapositivas de la asignatura de Electrónica Física, a priori no debería haber erratas ya que he corregido todas las que encontrábamos tanto yo como otros compañeros. En principio están todas las ecuaciones corregidas aunque siempre es bueno contrastar con las diapositivas de clase, también es recomendable añadir comentarios que se hayan hecho en clase ya que alguno se me ha podido escapar y son puntos que restan en el examen.\\\\
Las preguntas principales de teoría (las que han caído otros años) están marcadas con los títulos en \textcolor{red}{rojo} para que sea más fácil identificarlas. Que sean preguntas individuales no quiere decir que no puedan pedirse varias en una misma pregunta de teoría en el examen.\\\\
Los apuntes están divididos en Física de Semiconductores, que gran parte se ve en Física del Estado Sólido I; componentes electrónicos básicos, que es donde mayor cantidad de ejercicios diferentes hay; y dispositivos electrónicos especializados que únicamente es teoría y no hay ejercicios (suelen caer dos puntos de la parte de teoría del examen).\\\\

\begin{flushright}
    Mucha suerte a todos y espero que os sirvan por mucho tiempo.
\end{flushright}

\begin{flushright}
    Chinri
\end{flushright}

\vfill
\begin{flushleft}
        % Nota de derechos y autoría
    \noindent\textbf{Licencia y Derechos de Autor} \\
    \vspace{0.2cm}
    \small
    Este documento y su código fuente son de elaboración propia. 
    Se distribuyen bajo una licencia \\ \textbf{Creative Commons Attribution-NonCommercial-ShareAlike 4.0 International} (CC BY-NC-SA 4.0).
    
    \vspace{0.4cm}
    
    % Esto inserta los iconos típicos de Creative Commons de forma limpia usando texto
    \includegraphics[width=0.25\textwidth]{by-nc-sa.pdf}
    
    \vspace{0.4cm}
    \footnotesize
    Eres libre de compartir y adaptar este material siempre que se reconozca la autoría original, 
    no se utilice con fines comerciales y se distribuya bajo la misma licencia.
\end{flushleft}